%=====================================================================
\chapter{The Boundary With Other Courses}\label{app:boundary}
%=====================================================================
\normalfigures

This course sits early in the degree --- second semester --- and three later
courses depend on it. It therefore assumes little and promises a great deal,
and both are recorded here: what is assumed, what is delivered, and to whom,
so that a lecturer can see what may be taken for granted on each side and a
student can see why a familiar topic reappears.

\section{The Edition History}

\rowcolors{2}{white}{lightbg}
\begin{longtable}{@{}L{1.5cm}L{3.1cm}L{3.8cm}L{1.9cm}L{3.0cm}@{}}
\toprule
\rowcolor{primary!15}
\textbf{Edition} & \textbf{Status} & \textbf{Degrees} & \textbf{Semester}
& \textbf{Content published?} \\
\midrule
\endhead
\textbf{2017} & Computing Core, 3+1 & BSCS, BSSE, BSIT & IV, IV, \textbf{V}
  & yes \\
\textbf{2023} & Computing Core, 4 (3-3) & all ten computing BS, and ADCS &
  \textbf{2} & yes --- the 2017 text \\
2025 & Major Compulsory, 3+1 & BSCS and ADCS only --- \textbf{no BSIT} & III
  & \textbf{no} \\
\bottomrule
\end{longtable}

\smallskip
The department teaches CS-103 at \textbf{BSIT, second semester, 4 credit hours
(3+1)}, which matches HEC 2023. HEC 2017 placed it at BSIT semester V. The 2023
course content is the 2017 text word for word; 2025 removed it. The Course
Specification at the front records all six defects found in these outlines.

\begin{alertbox}[One booklet, two prerequisites]
HEC 2017's course description page gives \textbf{``Prerequisites: None''}.
The same booklet's scheme-of-studies tables, on every page that lists the
course, give \textbf{``Data Structures \& Algorithms''}. HEC 2023 gives
``None''.

\smallskip
At second semester Data Structures has not been taken, so this handout
\textbf{assumes no data-structures course}. The two chapters that would most
naturally lean on one are written to stand alone: \chref{indexes} explains a
B-tree by its fan-out and depth, with Appendix~B's logarithm, and never by an
algorithm; \chref{transactions} explains locks without operating-systems
vocabulary. Recorded here so that a reviewer comparing the two documents finds
the answer already written.
\end{alertbox}

\section{Programming Fundamentals --- the Real Prerequisite}

\begin{conceptbox}[The rule]
\textbf{Programming Fundamentals owns the program. This course owns what the
program keeps.}

\smallskip
A student arriving here should be able to write a loop, read a file, and call
a function. Nothing more is used: Python appears in the harnesses, which
students run and read rather than write, and in two laboratories ---
\chref{whynotfiles}'s file-based system and \chref{project}'s injection
demonstration --- where the programs are a dozen lines.
\end{conceptbox}

\section{Data Structures --- Not Assumed}

\rowcolors{2}{white}{lightbg}
\begin{longtable}{@{}L{5.4cm}L{3.4cm}L{5.6cm}@{}}
\toprule
\rowcolor{primary!15}
\textbf{Topic} & \textbf{Owned by} & \textbf{Here} \\
\midrule
\endhead
B-trees: insertion, splitting, balance & Data Structures & \emph{not taught};
  \chref{indexes} needs only fan-out and depth \\
hashing & Data Structures & mentioned, not used \\
sorting algorithms & Data Structures & \chref{grouping} measures what a sort
  \emph{costs} --- \sqlkw{TEMP B-TREE} --- not how it works \\
logarithms & Appendix~B & taught here, as ``how many halvings'' \\
\bottomrule
\end{longtable}

\smallskip
\textbf{The useful repetition:} a student who takes Data Structures later will
meet the B-tree's insertion algorithm, and will already know why it matters ---
\chref{indexes} measured a lookup costing 0.057\,ms in a million rows, and a
scan 199.90\,ms.

\section{E-Commerce Application Development --- Downstream}

\begin{conceptbox}[The rule, as that course states it]
\textbf{``Database Systems owns the model. This course owns the consequences of
somebody else's model.''}

\smallskip
E-Commerce Application Development's own Appendix~F says that, and lists what
it assumes: \emph{``SQL selects and joins, primary and foreign keys,
normalisation, what an index is.''} This course must deliver all four, and
does.
\end{conceptbox}

\rowcolors{2}{white}{lightbg}
\begin{longtable}{@{}L{4.6cm}L{4.6cm}L{5.2cm}@{}}
\toprule
\rowcolor{primary!15}
\textbf{E-Commerce assumes} & \textbf{Delivered in} & \textbf{And measured} \\
\midrule
\endhead
SQL selects and joins & \chref{select}, \chref{joins}, \chref{subqueries},
  \chref{grouping} & \sqlkw{NOT IN} returning 0 where \sqlkw{NOT EXISTS}
  returned 3 \\
primary and foreign keys & \chref{relational}, \chref{constraints} & foreign
  keys off by default; a \texttt{TEXT PRIMARY KEY} storing null \\
normalisation & \chref{fds} to \chref{decomposition} & a wide table
  162$\times$ slower on a lookup; 143{,}768 spurious tuples \\
what an index is & \chref{indexes} & the crossover --- which E-Commerce
  then finds again, in a product catalogue \\
\bottomrule
\end{longtable}

\smallskip
\textbf{The useful repetition}, from the downstream side: E-Commerce tells its
students that ``every Database Systems student has been told that indexes make
queries faster'', and makes that a measurement with a crossover in it. In this
course they have already measured it --- so the E-Commerce chapter becomes a
second instance of a known result, on a schema they did not design, rather than
a surprise.

\section{Information Security --- Later}

\begin{conceptbox}[The rule]
\textbf{Information Security owns the principles. This course owns one
database, and what can be done to it from inside.}

\smallskip
\chref{project} shows SQL injection against the student's own file, under the
scope statement of Appendix~A, and stops there. Threat modelling,
authentication, cryptography and the law are Information Security's; this
course shows the one defence that belongs to the person writing the query ---
the parameter --- and the access control a server adds that SQLite lacks.
\end{conceptbox}

\section{Software Engineering --- Later}

\begin{tight}
  \item \textbf{Assumed there, delivered here:} that a schema is designed
        before it is built (\chref{ermodel} to \chref{mapping}), and changed by
        numbered, reviewed migrations (\chref{evolution}).
  \item \textbf{Taught properly there, used informally here:} version
        control. \chref{evolution}'s first rule --- never edit a migration that
        has run anywhere --- is a version-control rule as much as a database
        one, and Software Engineering is where the tooling behind it is taught.
\end{tight}
